% ==============================================================================
% ================================ AUFGABE 5 ===================================
% ==============================================================================

%\chapter{Aufgabe 5: Flipflops – Schieberegister}
%\label{chap:aufgabe5}
%\thispagestyle{fancy}

\section{Aufgabe 5: Flipflops – Schieberegister}
In diesem Kapitel wird eine Schieberegisterkette aus drei D‑Flipflops konstruiert, die am seriellen Ausgang periodisch die Bitfolge \(010011\) liefert.  
Die Aufgabe gliedert sich in drei Teile: Erstellung der Zustandsfolgetabelle, Ermittlung des KV‑Diagramms für den Eingang des ersten D‑Flipflops und zeichnerische Darstellung der Digitalschaltung.
Die theoretischen Grundlagen zu dieser Aufgabe entstammen dem Studienheft MRT12 \cite{2}.
% ======================================= 5a =====================================
\subsection{Aufgabe 5a: Zustandsfolgetabelle}
Die gewünschte serielle Ausgangsfolge lautet: \hspace{4cm} \[Y = 0, 1, 0, 0, 1, 1, 0, 1, 0, 0, 1, 1, \dots\]
Die Schieberegisterkette besteht aus drei D‑Flipflops mit den Ausgängen \(Q_2, Q_1, Q_0\).  
Der serielle Ausgang ist \(Y = Q_2\).  
Der Zustand der Kette wird durch das Tupel (\(\mathit{Z}_2, \mathit{Z}_1, \mathit{Z}_0\)) beschrieben.

Die Zustandsfolge ergibt sich durch die drei aufeinanderfolgenden Bits der Ausgangsfolge.  
Für eine Periode der Länge 6 ergeben sich folgende Zustände:
\begin{table}[h]
\centering
\begin{tabular}{c|c|c|c|c}
\toprule
\textbf{Takt} & \textbf{Zustand \(Z_2 Z_1 Z_0\)} & \textbf{Ausgang \(Y = Z_2\)} & \textbf{nächster Zustand} & \textbf{Eingang \(D_0\)} \\
\midrule
0 & 010 & 0 & 100 & 0 \\
1 & 100 & 1 & 001 & 1 \\
2 & 001 & 0 & 011 & 1 \\
3 & 011 & 0 & 110 & 0 \\
4 & 110 & 1 & 101 & 1 \\
5 & 101 & 1 & 010 & 0 \\
\bottomrule
\end{tabular}
\caption{Zustandsfolgetabelle für die Schieberegisterkette (\textit{Quelle: Eigene Darstellung})}
\label{tab:schieberegister_zustaende}
\end{table}
Die Zustände wiederholen sich nach sechs Takten, sodass die Ausgangsfolge \(010011\) periodisch am seriellen Ausgang erscheint.
% ======================================= 5b =====================================
\subsection{\texorpdfstring{Aufgabe 5b: KV-Diagramm für den Eingang \(D_0\)}{Aufgabe 5b: KV-Diagramm für den Eingang D0}}
Der Eingang des ersten D‑Flipflops \(D_0\) wird aus dem aktuellen Zustand \((\mathit{Z}_2, \mathit{Z}_1, \mathit{Z}_0)\) gebildet.  
Aus der Zustandsfolgetabelle ergibt sich die folgende Wahrheitstabelle für \(D_0\):
\begin{table}[h]
\centering
\begin{tabular}{c|c|c|c}
\toprule
\(Z_2\) & \(Z_1\) & \(Z_0\) & \(D_0\) \\
\midrule
0 & 0 & 0 & 0 \\
0 & 0 & 1 & 1 \\
0 & 1 & 0 & 0 \\
0 & 1 & 1 & 0 \\
1 & 0 & 0 & 1 \\
1 & 0 & 1 & 0 \\
1 & 1 & 0 & 1 \\
1 & 1 & 1 & 0 \\
\bottomrule
\end{tabular}
\caption{Wahrheitstabelle für \(D_0\) (\textit{Quelle: Eigene Darstellung})}
\label{tab:d0_wahrheit}
\end{table}
Das zugehörige KV‑Diagramm für 3 Variablen (\textit{\(Z_2, Z_1, Z_0\)}) wird wie folgt aufgestellt.  
Die Anordnung der Spalten folgt dem Gray‑Code \(00, 01, 11, 10\) für \(Z_1 Z_0\) und die Zeilen unterscheiden \(Z_2\).
\begin{table}[h]
\centering
\begin{tabular}{c|cccc}
\multicolumn{1}{c}{} & \multicolumn{4}{c}{\(Z_1 Z_0\)} \\
\cline{2-5}
\(Z_2\) & 00 & 01 & 11 & 10 \\
\hline
0 & 0 & 1 & 0 & 0 \\
1 & 1 & 0 & 0 & 1 \\
\hline
\end{tabular}
\caption{KV-Diagramm für \(D_0\) (\textit{Quelle: Eigene Darstellung})}
\label{tab:d0_kv}
\end{table}
Die Einsen befinden sich bei:
- \(Z_2 Z_1 Z_0 = 001\)
- \(Z_2 Z_1 Z_0 = 100\)
- \(Z_2 Z_1 Z_0 = 110\)
Durch Zusammenfassen der beiden Einsen in der unteren rechten Ecke (100 und 110) ergibt sich der Block \(Z_2 \overline{Z_0}\).  
Die Eins bei 001 bleibt isoliert (\textit{\( \overline{Z_2} \overline{Z_1} Z_0 \)}).
Die minimierte Gleichung lautet somit: \hspace{4cm} \[\boxed{D_0 = Z_2 \cdot \overline{Z_0} \;\cup\; \overline{Z_2} \cdot \overline{Z_1} \cdot Z_0}\]
% ======================================= 5c =====================================
\subsection{Aufgabe 5c: Digitalschaltung}
Die Schaltung besteht aus drei D‑Flipflops, die in Reihe geschaltet sind.  
Der Ausgang \(Q_2\) des letzten Flipflops stellt den seriellen Ausgang \(Y\) dar.  
Der Eingang des ersten Flipflops \(D_0\) wird durch die minimierte Logik aus den Zuständen \(Z_2, Z_1, Z_0\) erzeugt.
\begin{figure}[h]
\centering
\includegraphics[width=0.75\textwidth]{asset/image/Aufgabe5.png}
\caption{Schaltbild der Schieberegisterkette mit drei D-Flipflops und der Logik für \(D_0\) (\textit{Quelle: Eigene Darstellung})}
\label{fig:schieberegister_schaltung}
\end{figure}
Die Schaltung erzeugt bei korrekter Verkabelung die periodische Ausgangsfolge \(010011\) am seriellen Ausgang \(Y\).

% ==============================================================================
% ================================ AUFGABE 5 ===================================
% ==============================================================================